\begin{center}
    ABSTRACT
\vspace{5mm} %5mm vertical space
\end{center}

Post-harvest mould caused by \textit{Botrytis cinerea} is a major source of food loss in fresh produce supply chains. This thesis investigates the energy cost of TinyML framework execution on an ESP32-based autonomous mould detection node operating without cloud connectivity. A three-node wireless sensing mesh --- each node integrating a DHT22, SGP30 VOC sensor, and MQ3 ethanol sensor on a custom PCB --- was used to collect environmental data across five controlled strawberry storage experiments. Three frameworks were evaluated: TF Lite Micro (inference only), TinyOL (output-layer adaptation), and AIfES (full on-device training and inference).

Energy was measured using a Nordic Semiconductor PPK2 power analyser. AIfES inference consumed \SI{6.3}{\micro\joule} per forward pass, outperforming INT8-quantised TF Lite Micro (\SI{8.9}{\micro\joule}) due to the ESP32 hardware FPU. Full AIfES on-device training consumed \SI{194.8}{\micro\joule} per gradient step --- 107$\times$ more than a TinyOL update --- but less than 0.06\% of total daily node energy. Classification accuracy on a held-out test set reached 94.0\% for AIfES inference, 93.7\% for TF Lite Micro, 86.1\% for TinyOL, and 77.1\% for full on-device training. Battery lifetime is determined entirely by sensing hardware: the MQ3 heater alone accounts for 61.5\% of node current, limiting a 10,000\,mAh battery to approximately 16~hours. Replacing the MQ3 with a MEMS sensor and duty-gating the SGP30 extends this to 4.2~days with no PCB redesign.

\vspace{5mm}
\noindent {\bf Keywords:} TinyML, edge AI, energy profiling, ESP32, mould detection, AIfES, TinyOL, TF Lite Micro, post-harvest monitoring.
